% ---------------------------------------------------------------------
% methodology.tex -- the extended genome, the rediscovery argument, and
% the hardware layer.  \input at the end of 3-methods.tex.
% Style rule: never use the em-dash.
% ---------------------------------------------------------------------

\subsection{An exact-path lasso solver as a search-space example}
\label{sec:method:rediscover}

Reachability is not rediscovery: the expert's decisions must be expressible
as mutations, and the fitness must make their effects visible. An exact-path
lasso solver illustrates this design. Each of the four decisions in
Table~\ref{tab:rediscover} becomes a mutation of one layer.

\begin{table}[htbp]\centering\footnotesize
\caption{The four decisive steps of an expert derivation, each reachable
as a single typed mutation. D3 and D4 are visible to any wall-clock
fitness; D5 and D6 are not, which is the reason the fitness changes.}
\label{tab:rediscover}
\begin{tabular}{@{}p{1.42cm}p{1.24cm}p{3.05cm}p{3.30cm}p{3.04cm}@{}}
\toprule
Step & Layer patched & The mutation & Gate rule that keeps it honest &
Why the fitness sees it \\
\midrule
D3 carry the factorization &
3 state &
\texttt{factorization.\allowbreak update}: refactorize to insert\_\allowbreak delete &
drift check on the maintained factor &
per-step cost, visible immediately \\
\addlinespace
D4 free screening bound &
4 restrict &
\texttt{rule}: none to equiangular\_\allowbreak bound &
admitted only if provably safe, or paired with a full re-admission pass,
checked statically &
step-count identity proves exactness, so the win is pure \\
\addlinespace
D5 hybridize &
5 handoff &
\texttt{handoff}: null to \{to: working\_\allowbreak set, predicate\} &
the existing switch\_k rule generalized: exact prefix exempt, tail passes
the full gate &
neither engine dominates, so only a sweep shows it \\
\addlinespace
D6 measure, do not bound &
5 handoff &
\texttt{predicate.\allowbreak source}: bound to counted &
none needed; it is a source swap &
only oracle regret on a sweep sees it \\
\bottomrule
\end{tabular}
\end{table}

The D6 decision illustrates why fitness matters. The gene value
$\texttt{source} \in \{\text{bound}, \text{fitted}, \text{probe},
\text{counted}\}$ is an enumeration, so swapping it is a legal single
mutation. Under a wall-clock fitness, the negated shifted geometric mean
of time-to-certificate over a training set, \texttt{bound} looks fine,
because it is calibrated on the instance where the bound is tight. Under
regret against a per-problem oracle across $54$ shape-and-grid
combinations at two thread counts, \texttt{bound} costs $14\times$ to
$19\times$ and \texttt{counted} costs $1.10\times$ with a worst case of
$1.64\times$. The same mutation gets opposite verdicts under the two
fitness measures; the chosen fitness thus determines which predicate a search
would favor.

The example shows how an expert's four-step derivation can be expressed as
four single mutations; it does not show that a search run rediscovered them.
An ablation of that discovery claim would withhold the four decisive gene-pool
cards, compare LLM mutation without a gene pool, and run a random proposer
over the same typed space. SimpleTES \citep{Ye2026StructuredScaling} is another
algorithm-discovery system with a lasso-path solver and offers an external
comparison target.

\subsection{Where hardware kernel fusion comes in}
\label{sec:method:hardware}

Fusion is layer~8, and layer~8 is not a post-processing step applied to
the winning algorithm. It is one of three hardware genes, all set from the
same quantity, the work per step of whatever inner loop layers~1
through~6 produced.

\begin{table}[htbp]\centering\footnotesize
\caption{The three hardware genes. Each has a GPU form and a CPU form, and
each has a measured cliff, which is what distinguishes a gene from an
environment default.}
\label{tab:hwgenes}
\begin{tabular}{@{}p{1.42cm}p{1.95cm}p{3.35cm}p{2.35cm}p{2.95cm}@{}}
\toprule
Gene & Question & GPU form & CPU form & The cliff that proves it is a gene \\
\midrule
7 admission &
cross the device at all? &
streamed bytes against the bandwidth difference, decided before touching
the runtime &
stay &
importing the runtime costs more than a small solve \\
\addlinespace
8 width &
how much work per launch or barrier? &
fusion tile: $T$ iterations per launch on a $(B{+}2T)^2$ tile with halo in
shared memory &
BLAS thread count from work per step &
the wrong device's tile costs $1.06\times$ to $1.63\times$; $64$ threads
instead of $16$ costs $103\times$ \\
\addlinespace
6 contract &
what precision, and where &
per-expression fastmath scope &
same &
fastmath on a Schur complement turned a $10^{-14}$ certificate into
$10^{-2}$ \\
\bottomrule
\end{tabular}
\end{table}

The fusion tile and the thread count are the same gene at two
granularities. Both answer how wide the parallel region should be relative
to the work the algorithm gives it, both are read off the device profile,
and both have a cliff when set by an environment default instead of by the
algorithm. That is what makes the RTX~4090 against H100 tile result and
the $16$ against $64$ thread result one finding rather than two anecdotes.

Algorithm~\ref{alg:optevolve} closes the loop by feeding the measured
$t(\mathrm{operator},\mathrm{device})$ back into the layer-5 routing
predicate. A better fused tile lowers $t$ for an iterative candidate and
moves the crossover at which it beats another algorithm. Hardware choice
therefore changes the route, not just the cost of a fixed solver
(Equation~\ref{eq:costmodel}).
